% ===========================================================================
\chapter{A pretend model}\label{ch:model}

To the agent, a model is something that takes the conversation so far and
gives back the next reply. How it does that is its own business: Claude
does it over the network, but a script can do it too, and for building and
testing the agent a script is far better. It costs nothing, it answers
instantly, and it says exactly what we want it to say, including the odd
things a real model only does once in a while, like asking for a tool that
doesn't exist.

\section{What a model is given}

A request to the model holds three things: the \emph{system prompt}, the
standing instructions that say who the model is in this conversation and
how to work; the conversation; and the tools it may ask for, each
described as a JSON object (Chapter~\ref{ch:tools}). A request only makes
sense when it is the model's turn to speak.

\begin{newfile}{model.nex}
-- model.nex: the model, seen from the agent: something that takes the
-- conversation so far and gives the next reply.
--
-- Claude_Model (claude_model.nex) asks Claude over the network. The
-- models here don't: Scripted_Model plays a part written by hand,
-- Replay_Model plays back a session recorded earlier, and
-- Recording_Model records a session while another model does the work.

intern Messages

-- Everything the model is given for one reply.
class Request
create
  make(system: String, conversation: Conversation, tools: Array[Map[String, Any]])
  require
    models_turn: conversation.models_turn()
  do
    this.system := system
    this.conversation := conversation
    this.tools := tools
  end

feature
  system: String                    -- the standing instructions
  conversation: Conversation
  tools: Array[Map[String, Any]]    -- the tools it may ask for
end
\end{newfile}

\section{The Model class}

\code{Model} is a \emph{deferred} class: it says what every model can do,
without saying how. Its \code{reply} routine has no body (the word
\code{deferred} stands in for one); each kind of model will provide its
own. A class that \code{inherit}s \code{Model} and provides every deferred
routine is a model the agent can use, whatever it is underneath.

\begin{append}{model.nex}
deferred class Model
feature
  -- The model's next reply.
  reply(request: Request): Reply
  require
    ready: can_reply()
  deferred

  -- Whether it can give another reply. A model over the network always
  -- can (though the network may fail); a script can run out.
  can_reply(): Boolean do
    result := true
  end

  -- Who is answering, for the transcript: "claude-opus-5-5".
  name(): String deferred
end
\end{append}

The precondition on \code{reply} deserves a closer look. A script can run
out of lines, so a scripted model must not be asked for a reply it doesn't
have. The obvious place to say that is on the scripted model's own
\code{reply}. But Nex won't allow it, and rightly so: code that holds a
\code{Model} (the agent) knows only \code{Model}'s contract, so a
precondition that a particular kind of model adds would be a rule the
caller can't see. An heir may only accept \emph{more} than its parent, never
less. So the rule goes on \code{Model} itself, in a form every model can
fill in: \code{reply} may be called when \code{can_reply()} is true. Each
model says when that is, and the agent asks before every reply.

\section{A model that reads a script}

A \code{Scripted_Model} is given its replies in advance, and gives them
back one at a time, whatever it is asked: an actor reading a script. It
remembers what it was sent each time, as JSON text, so that our checks can
look at what the agent told the model.

What is sent each time is the whole conversation, and all but its last
message were sent before. So the interesting part is the last user
message: the task, the first time; then the results of the tools the model
asked for. \code{last_message_json} gets it as JSON text. (It turns the
message into JSON text, reads it back and turns it into text again. That
looks pointless, but it means two equal messages always give equal text,
however they were built, which is what we need to compare them.)

\begin{append}{model.nex}
-- The JSON of the message the user sent last in `request`, the part of
-- the request that is new. JSON text read back is used for both sides of
-- a comparison, so that two equal messages give equal text.
function last_message_json(request: Request): String do
  result := json_text(json_value(json_text(request.conversation.last_user_message())).raw())
end

-- A model that gives the replies it was given, in order, whatever it is
-- asked: an actor reading a script. It keeps what it was sent each
-- time, so a check can look at it.
class Scripted_Model inherit Model
create
  make(replies: Array[Reply]) do
    this.replies := replies
  end

feature
  replies: Array[Reply]
  sent: Array[String]   -- the new message of each request, as JSON

  name(): String do
    result := "a script"
  end

  can_reply(): Boolean do
    result := sent.length < replies.length
  end

  reply(request: Request): Reply do
    result := replies.get(sent.length)
    sent.add(last_message_json(request))
  end
end
\end{append}

\code{Scripted_Model} gives \code{can_reply} a new body: an heir may
\emph{redefine} a routine it inherits, as long as it keeps its signature.
And because the agent always asks \code{can_reply} before \code{reply},
\code{replies.get(sent.length)} is always in range.

Here is a two-line script. We ask for a reply, answer its tool call by
hand, and ask again:

\begin{tryit}{try_script.nex}
intern Model

let model: Scripted_Model := create Scripted_Model.make([
  use_tool("Let me look at the loans.", "read_file", '{"path": "loans.nex"}'),
  say("I have seen enough.")
])
let c: Conversation := create Conversation.make()
c.add_user_text("What does loans.nex do?")

let first: Reply := model.reply(create Request.make("Be brief.", c, []))
console.print_line(first.text() + " (" + first.stop_reason() + ")")
c.add_reply(first)
c.add_tool_results([create Tool_Result.make(first.tool_calls().get(0).id,
  "class Loan ...", false)])

let second: Reply := model.reply(create Request.make("Be brief.", c, []))
console.print_line(second.text() + " (" + second.stop_reason() + ")")
console.print_line("can reply again: " + model.can_reply().to_string)
across model.sent as message do
  console.print_line("sent: " + message)
end
\end{tryit}
\tryoutput{try_script}

The model said what the script said, and kept what it was sent: first the
question, then the result of the tool. In Chapter~\ref{ch:loop} the agent
will do the part we just did by hand.
